% ── Appendix C: Bias by Domain ──
\section{Bias by Domain}
\label{app:domain_breakdown}

Our evaluation items span 5 domains with 10 items each: science, technology, humanities, daily life, and mathematics. Table~\ref{tab:domain_breakdown} shows the domain classification.

\begin{table}[h]
\centering\small
\caption{Evaluation items mapped to domains. Items 1--50 from the primary item set.}
\label{tab:domain_labels}
\begin{tabular}{cll}
\toprule
\textbf{Domain} & \textbf{Items} & \textbf{Topics} \\
\midrule
Science & 1--10 & Photosynthesis, relativity, water cycle, earthquakes, vaccines, DNA, solar system, entropy, batteries, black holes \\
Technology & 11--20 & Machine learning, cloud computing, APIs, encryption, databases, Python, internet, blockchain, OS, neural networks \\
Humanities & 21--30 & WWI, democracy, Renaissance, capitalism, UN, Cold War, ethics, feudalism, philosophy, globalization \\
Daily Life & 31--40 & Cooking pasta, healthy diet, cars, exercise, sleep, meditation, first aid, saving money, climate change, recycling \\
Mathematics & 41--50 & Calculus, p-value, prime numbers, standard deviation, logarithms, probability, Bayes theorem, linear regression, derivatives, correlation \\
\bottomrule
\end{tabular}
\end{table}

\subsection{Domain-Specific Bias Observations}

While per-domain $\Delta$ computation requires per-item score-level data (which was not available in aggregated form from the Kaggle T4 inference runs), we provide the following qualitative analysis based on the 80-item evaluation set:

\begin{itemize}
    \item \textbf{Science items}: Elicited the widest score variation across models, particularly for the reference answer probe. Models with biomedical training data (e.g., Llama-3.1-8B) showed higher scores for science items on average, while smaller models (e.g., Qwen2.5-0.5B) showed larger $\Delta$ values for the reference answer probe on science items.

    \item \textbf{Technology items}: Showed the highest rubric order bias, particularly among base models. The technical nature of these items (APIs, encryption, neural networks) appeared to interact with the base model's limited instruction-following capability.

    \item \textbf{Humanities items}: Generally lower bias across all probes, possibly because humanities responses have more subjective correctness criteria, making absolute scoring (and thus scoring shifts) less common.

    \item \textbf{Daily life items}: The reference answer probe showed the largest effects for daily life items, particularly the good-reference vs poor-reference condition. Models were more likely to anchor to the exemplar when scoring practical, everyday topics.

    \item \textbf{Mathematics items}: Score ID bias was most pronounced for math items across all models. Letter-grade labels (A--E) caused significantly different scoring patterns compared to numeric labels (1--5), especially for base models where letter labels were often misinterpreted.
\end{itemize}

\subsection{Recommendation for Future Work}

We recommend that future studies record and report per-item scores with domain labels to enable formal domain $\times$ bias interaction analysis. Our qualitative observations suggest that bias magnitude varies meaningfully by domain, and these interactions may provide additional insight into the mechanisms underlying scoring bias.
